%==============================================================================
%  Supplementary material for
%  "Three Mechanisms Inflate Reported Performance on a Widely Used Public
%   Chronic Kidney Disease Benchmark"
%
%  Build:  pdflatex supplement && pdflatex supplement
%
%  Every table here is generated from a committed CSV in reports/tables/;
%  the source file is named in a comment above each.
%==============================================================================
\documentclass[conference]{IEEEtran}

\usepackage{cite}
\usepackage{amsmath,amssymb}
\usepackage{graphicx}
\usepackage{booktabs}
\usepackage{longtable}
\usepackage{xcolor}
\usepackage{url}
\usepackage[hidelinks]{hyperref}

\graphicspath{{figures/}}
\newcommand{\tabhead}[1]{\textbf{#1}}
\renewcommand{\arraystretch}{1.15}

% Number supplementary floats S1, S2, ...
\renewcommand{\thetable}{S\arabic{table}}
\renewcommand{\thefigure}{S\arabic{figure}}
\renewcommand{\thesection}{S\arabic{section}}

\begin{document}

\title{Supplementary Material\\
Three Mechanisms Inflate Reported Performance\\
on a Widely Used Public Chronic Kidney Disease Benchmark}

% TODO(author): mirror the main manuscript's author block before submission.
\author{\IEEEauthorblockN{Author Name}
\IEEEauthorblockA{\textit{Department} \\
\textit{Institution}\\
City, Country \\
email@institution.edu}}

\maketitle

\section{Scope of this supplement}

This document holds material that supports the main text but would displace
its argument: the full feature registry, the complete provenance sensitivity
suite, candidate flagging values for the audit framework, the extended
discussion of synthetic-data pipelines, reproducibility detail, and the
reporting checklist.

Every quantity is generated from a committed result file. Where an analysis
could not be completed, this supplement says so rather than omitting the
question.

%==============================================================================
\section{Feature-availability and leakage registry}

The registry (\texttt{data/feature\_registry.csv}) records 18 fields for each
of the 29 columns. Table~\ref{tab:sregistry} summarises the classification;
the full file, including per-variable evidence and uncertainty notes, ships
with the code.

% Source: reports/tables/table_47_registry_summary.csv
\begin{table}[h]
\caption{Variables per leakage category.}
\label{tab:sregistry}
\centering
\footnotesize
% Generated by scripts/24_make_latex_tables.py -- do not edit.
\begin{tabular}{lr}
\toprule
\tabhead{Leakage category} & \tabhead{Variables} \\
\midrule
Direct outcome copy & 2 \\
Deterministic derivative & 2 \\
KDIGO criterion input & 2 \\
Urine sediment (possible criterion) & 4 \\
Correlated renal biomarker & 4 \\
Consequence of advanced disease & 4 \\
Uncertain provenance & 6 \\
Legitimate pre-index & 5 \\
\bottomrule
\end{tabular}

\end{table}

Feature configurations are derived from this file rather than maintained
separately. \texttt{tests/test\_registry.py} asserts that the derivation
reproduces the configurations used in the reference analysis exactly; if a
classification is edited without updating the configurations, or vice versa,
the test suite fails.

%==============================================================================
\section{Provenance sensitivity: full results}

Section~V-D of the main text reports the headline sensitivity analyses.
Table~\ref{tab:loo} gives the complete leave-one-variable-out results,
generated from the result file rather than transcribed.

% Generated by scripts/24_make_latex_tables.py -- do not edit.
% Source: reports/tables/table_44_provenance_sensitivity.csv
\begin{table}[t]
\caption{Leave-one-variable-out matching, complete results. Each row removes one containment variable and re-runs the match. The match fraction is 1.000 in every case; \emph{unique pins} is the discriminating statistic. Ordered by unique pins, fewest first.}
\label{tab:loo}
\centering
\footnotesize
\begin{tabular}{lccc}
\toprule
\tabhead{Variable removed} & \tabhead{$k$} & \tabhead{Match} & \tabhead{Unique pins} \\
\midrule
\texttt{age} & 12 & 1.000 & 133 \\
\texttt{bgr} & 12 & 1.000 & 159 \\
\texttt{bu} & 12 & 1.000 & 170 \\
\texttt{al} & 12 & 1.000 & 171 \\
\texttt{sod} & 12 & 1.000 & 172 \\
\texttt{sg} & 12 & 1.000 & 179 \\
\texttt{hemo} & 12 & 1.000 & 179 \\
\texttt{su} & 12 & 1.000 & 180 \\
\texttt{sc} & 12 & 1.000 & 181 \\
\texttt{pcv} & 12 & 1.000 & 184 \\
\texttt{rbcc} & 12 & 1.000 & 184 \\
\texttt{wbcc} & 12 & 1.000 & 185 \\
\texttt{pot} & 12 & 1.000 & 187 \\
\bottomrule
\end{tabular}
\end{table}


\subsection{Matched-pair audit}

\texttt{table\_45\_matched\_pair\_audit.csv} contains twelve reproducibly
sampled matched pairs (seed fixed in the script), each showing the analysed
patient's published interval alongside the candidate record's value for every
matching variable. Readers can regenerate and inspect it with
\texttt{python scripts/21\_provenance\_sensitivity.py}.

\subsection{A null that cannot be used}

Permuting whole candidate rows does not provide a null for this statistic:
maximum bipartite matching depends only on the multiset of candidate records,
not on their order, so the statistic is invariant by construction. We
implemented it, observed exactly zero variance, and replaced it with the
copula null reported in the main text. The check is retained as
\texttt{row\_permutation\_is\_degenerate} in
\texttt{provenance\_sensitivity.py}, so the claim is verifiable rather than
asserted. It is reported here, and not in the main text, because it is a
property of the estimator rather than a competing scientific null.

\subsection{Which nulls are implemented, and whether they work}

Two nulls are implemented and reported. \emph{Independent column permutation}
shuffles each candidate column separately, destroying every cross-variable
relationship; beating it establishes only that candidate records are
internally coherent. The \emph{Gaussian-copula null} is the stricter one:
within each outcome stratum the 13 containment variables are rank-transformed
to normal scores, a correlation matrix is estimated on those scores, synthetic
normal draws are generated through its Cholesky factor, and each column is
mapped back through its own empirical marginal using nearest-value
quantiles, so that every synthetic value is a value the column actually
takes. Missingness is preserved per column and per record. Seeds derive from
the master seed; 200 draws are used.

Nearest-value mapping matters. Linear interpolation would place synthetic
values \emph{between} the release's discrete bin representatives, where no
published interval can contain them; the null would then be depressed for a
purely mechanical reason and would flatter the observed result. With the
discrete support preserved, the null rises slightly, which is the
conservative direction.

A null is only as good as its construction, so we measured whether it
preserves what it claims (Table~\ref{tab:nulldiag}). Across the 78 variable
pairs the mean absolute difference between the real and synthetic Spearman
correlations is 0.031, against real correlations spanning $-0.73$ to $+0.87$;
the largest single discrepancy is 0.168. No synthetic value falls outside its
column's observed support. The null therefore reproduces the cohort's shape
closely while destroying record identity, which is what makes exceeding it
informative.

% Generated by scripts/24_make_latex_tables.py -- do not edit.
% Sources: table_53_copula_null_diagnostics.csv,
%          table_46_provenance_null_summary.csv
\begin{table}[h]
\caption{Null models actually implemented, and whether the dependence-preserving null preserves what it claims. Row-order permutation is excluded: bipartite matching is invariant to candidate row order, so it is a property of the estimator, not a null (Section~S3-B).}
\label{tab:nulldiag}
\centering
\footnotesize
\setlength{\tabcolsep}{2pt}
\begin{tabular}{@{}p{4.1cm}cr@{}}
\toprule
\tabhead{Quantity} & \tabhead{Value} & \tabhead{Basis} \\
\midrule
Independent column permutation & 0.003 $\pm$ 0.004 & 200 draws \\
\quad exceedances $\geq$ observed & 0 & $p \leq 0.005$ \\
Gaussian-copula null & 0.059 $\pm$ 0.014 & 200 draws \\
\quad exceedances $\geq$ observed & 0 & $p \leq 0.005$ \\
\midrule
Variable pairs compared & 78 & Spearman \\
Mean $|$real $-$ null$|$ correlation & 0.031 & retained \\
Largest single discrepancy & 0.168 & \texttt{su~bgr} \\
Synthetic values off observed support & 0 & exact \\
\bottomrule
\end{tabular}
\end{table}


%==============================================================================
\section{All registry-derived configurations}

Table~\ref{tab:restrictedfull} reports every configuration derived from the
corrected registry, with the full metric set. The main text shows seven of
these; the remaining rows separate the two limbs of the KDIGO definition so
that a reader can see each removed on its own. \emph{Saturated} marks a
configuration whose lower interval bound is at or above 0.99. The two
unsaturated rows are the ones restricted to information plausibly available
before the renal work-up.

% Generated by scripts/24_make_latex_tables.py -- do not edit.
% Source: reports/tables/table_49_restricted_featuresets.csv
\begin{table*}[t]
\caption{Every registry-derived configuration, best model per row. 200 unique patients (128 cases); 95\% stratified patient-level bootstrap intervals. \emph{Saturated} marks a configuration whose lower interval bound is at or above 0.99, i.e.\ one retaining no measurable headroom on this sample.}
\label{tab:restrictedfull}
\centering
\footnotesize
\setlength{\tabcolsep}{4pt}
\begin{tabular}{p{4.4cm}cl ccccc c}
\toprule
\tabhead{Feature set} & \tabhead{$k$} & \tabhead{Model} & \tabhead{ROC-AUC (95\% CI)} & \tabhead{PR-AUC} & \tabhead{Sens.} & \tabhead{Spec.} & \tabhead{Brier} & \tabhead{Sat.} \\
\midrule
full\_valid (reference) & 25 & SVM & 1.000 (1.000--1.000) & 1.000 & 1.000 & 1.000 & 0.004 & yes \\
minus albuminuria information & 24 & SVM & 1.000 (1.000--1.000) & 1.000 & 1.000 & 1.000 & 0.004 & yes \\
minus criterion inputs only (strict) & 23 & SVM & 1.000 (1.000--1.000) & 1.000 & 1.000 & 1.000 & 0.004 & yes \\
minus all incorporation-risk variables & 19 & SVM & 1.000 (1.000--1.000) & 1.000 & 1.000 & 1.000 & 0.004 & yes \\
minus creatinine/eGFR information & 24 & SVM & 1.000 (1.000--1.000) & 1.000 & 1.000 & 1.000 & 0.003 & yes \\
minus consequences of advanced disease & 21 & SVM & 0.998 (0.995--1.000) & 0.999 & 0.969 & 0.986 & 0.017 & yes \\
minus incorporation and consequences & 15 & SVM & 0.994 (0.987--1.000) & 0.997 & 0.969 & 0.972 & 0.023 & no \\
conservative pre-index (no uncertain) & 5 & RF & 0.930 (0.894--0.962) & 0.968 & 0.852 & 0.875 & 0.095 & no \\
pre-index available only & 5 & RF & 0.930 (0.894--0.962) & 0.968 & 0.852 & 0.875 & 0.095 & no \\
minimal low-cost conservative & 3 & RF & 0.881 (0.833--0.923) & 0.944 & 0.773 & 0.764 & 0.124 & no \\
\bottomrule
\end{tabular}
\end{table*}


%==============================================================================
\section{Case-mix re-weighting}

Table~\ref{tab:casemixsim} reports measured performance under target severity
distributions, with the fitted models held fixed and the out-of-fold
predictions re-weighted. It is a resampling exercise and not prospective
evidence: it answers what would have been \emph{measured} on a differently
composed cohort using this model, not what a model trained on that cohort
would achieve. The last column is the honest limit on it --- with 9 stage-1
and 12 stage-2 cases in the benchmark, an early-heavy target reuses a few
patients many times, and the Monte Carlo intervals do not reflect that.

% Generated by scripts/24_make_latex_tables.py -- do not edit.
% Source: reports/tables/table_54_casemix_simulation.csv
\begin{table*}[t]
\caption{Measured performance under target severity distributions, models held fixed and out-of-fold predictions re-weighted. \textbf{Exploratory resampling, not prospective evidence.} Intervals are Monte Carlo only and do not reflect the thinness of the support: the benchmark holds 9 stage-1 and 12 stage-2 cases, so early-heavy targets reuse a few patients many times, as the last two columns show.}
\label{tab:casemixsim}
\centering
\footnotesize
\setlength{\tabcolsep}{4pt}
\begin{tabular}{p{2.6cm}p{3.3cm}cccc}
\toprule
\tabhead{Configuration} & \tabhead{Target mix} & \tabhead{Early frac.} & \tabhead{ROC-AUC (MC 95\%)} & \tabhead{Sens.} & \tabhead{Distinct / max reuse} \\
\midrule
low\_cost & as sampled (benchmark) & 0.16 & 0.993 (0.984--1.000) & 0.968 & 81 / 8$\times$ \\
low\_cost & community screening series & 0.68 & 0.989 (0.981--0.996) & 0.938 & 43 / 16$\times$ \\
low\_cost & early-heavy & 0.80 & 0.987 (0.979--0.994) & 0.923 & 38 / 16$\times$ \\
low\_cost & uniform across stages & 0.40 & 0.990 (0.981--0.997) & 0.950 & 75 / 10$\times$ \\
low\_cost & advanced-heavy & 0.00 & 0.994 (0.985--1.000) & 0.979 & 74 / 8$\times$ \\
laboratory & as sampled (benchmark) & 0.16 & 0.994 (0.986--0.999) & 0.969 & 81 / 8$\times$ \\
laboratory & community screening series & 0.68 & 0.985 (0.977--0.992) & 0.912 & 43 / 16$\times$ \\
laboratory & early-heavy & 0.80 & 0.984 (0.976--0.990) & 0.879 & 38 / 16$\times$ \\
laboratory & uniform across stages & 0.40 & 0.990 (0.983--0.996) & 0.934 & 75 / 10$\times$ \\
laboratory & advanced-heavy & 0.00 & 0.997 (0.990--1.000) & 0.991 & 74 / 8$\times$ \\
full\_valid & as sampled (benchmark) & 0.16 & 1.000 (1.000--1.000) & 1.000 & 81 / 8$\times$ \\
full\_valid & community screening series & 0.68 & 1.000 (1.000--1.000) & 1.000 & 43 / 16$\times$ \\
full\_valid & early-heavy & 0.80 & 1.000 (1.000--1.000) & 1.000 & 38 / 16$\times$ \\
full\_valid & uniform across stages & 0.40 & 1.000 (1.000--1.000) & 1.000 & 75 / 10$\times$ \\
full\_valid & advanced-heavy & 0.00 & 1.000 (1.000--1.000) & 1.000 & 74 / 8$\times$ \\
\bottomrule
\end{tabular}
\end{table*}


%==============================================================================
\section{Audit framework: candidate flagging values}

The main text reports the three measurements as quantities. Candidate
flagging values are given here because they are calibrated on a single
benchmark and should not be read as criteria:

\begin{itemize}
  \item multivariable lift below 0.05,
  \item 90\% of achievable gain reached below 25\% of the training data,
  \item standardisation shift above 0.05.
\end{itemize}

Applied to this benchmark, at least one value is exceeded for every
configuration. We repeat that this demonstrates internal consistency, not
external validity: the values have not been shown to separate informative
from uninformative benchmarks, and doing so requires a multi-dataset study
with known ground truth that we have not performed.

\subsection{The negative result, in full}

Standardising ROC-AUC to a published community screening case mix moves it by
at most 0.008 on this dataset --- it would not flag a benchmark that Section
V-B shows to be case-mix sensitive at a threshold. Standardising sensitivity
at a fixed operating point moves it by up to 0.056. The explanation is
specific to this sample: early-stage cases here have subgroup ROC-AUC 0.983
(21 cases), so re-weighting toward them changes few rankings. This is
\emph{not} evidence that AUC is generally insensitive to spectrum effects; a
case mix that pushes cases below controls does move AUC. The practical
implication is narrower: standardising discrimination alone can give false
reassurance, so an operating-point metric should be standardised alongside
it.

%==============================================================================
\section{Extension to synthetic tabular and synthetic patient data}

The failure modes documented in the main text are not specific to real data,
and the safeguards generalise to pipelines that generate or evaluate
synthetic tabular records. Synthetic-data evaluation is harder in one
respect: a generator trained on a source table can reproduce individual rows,
so ``is the evaluation set independent of the training set?'' becomes a
question about memorisation as well as provenance.

\textbf{Already implemented here}, in a form adapted to real-data provenance:
cryptographic dataset fingerprints (SHA-256 on every load, re-verified in
continuous integration); exact-row and near-duplicate detection across
training, validation, test and reference sets (the interval-containment
matching, a tolerance-based near-duplicate test calibrated against a
permutation null); entity-level overlap checks with held-out variable
agreement; feature--target leakage scans and detection of target proxies and
post-outcome variables (the registry and its guards); enforcement that all
preprocessing is fitted inside cross-validation folds; provenance manifests
recording origin, version, transformations and split history; and a
fail-closed gate that blocks evaluation when a cohort is not verified
independent.

\textbf{Not implemented here, and required for synthetic data:}
membership-inference and memorisation checks, which ask whether a synthetic
record discloses that a particular real patient was in the training set; and
nearest-neighbour similarity or distance-to-closest-record diagnostics, which
ask whether synthetic rows are near-copies of real ones. The latter is both a
privacy question and an evaluation-validity question --- a synthetic test set
whose records sit arbitrarily close to training records reproduces, in
subtler form, exactly the non-independent evaluation this paper documents.
The present work provides the provenance half of the problem, not the
memorisation half.

A caution carries over with the recipe. The most consequential mechanism in
this study was case composition, which no fingerprint, hash or duplicate scan
would detect. A synthetic cohort can be free of memorisation, pass every
overlap check, and still be uninformative because its severity distribution
makes the task trivial. Provenance gating and informativeness diagnostics are
therefore complementary rather than alternatives.

%==============================================================================
\section{Reproducibility}

\subsection{One-command reproduction}

\begin{verbatim}
python reproduce.py
\end{verbatim}

validates input hashes, preprocesses, runs the analysis stages and test
suite, regenerates tables and figures, builds the manuscript, and verifies
that generated results match manuscript values. Individual targets are
documented in \texttt{README.md}. A \texttt{Makefile} offers the same
targets for POSIX users, but \texttt{make} is absent from the reference
machine, so it is provided \emph{unverified}: \texttt{reproduce.py} is the
path we have run.

\subsection{Environment}

The reference environment is CPython 3.14.5 on Windows 11 with numpy 2.5.2,
pandas 3.0.5, scikit-learn 1.9.0 and scipy 1.18.0; the complete freeze is in
\texttt{requirements.lock.txt}, which also records the BLAS build.

Identical seeds reproduce identical predictions within an environment. Across
a rebuild from the same pinned versions, agreement is to about $10^{-16}$
rather than bit-for-bit, because floating-point summation order depends on
the installed BLAS. The regression test asserts $10^{-12}$ rather than
claiming exactness the project cannot deliver.

\subsection{Data access}

The two UCI releases are openly licensed (CC BY 4.0) and are retrieved by
scripts with recorded checksums. The MIMIC-IV Clinical Database Demo is
openly available from PhysioNet. One registered dataset --- a community-based
CKD screening series --- is access-restricted and was not obtained; the
manuscript states where that limits the work.

%==============================================================================
\section{Reporting checklist}

A completed TRIPOD+AI checklist ships as
\texttt{reports/tables/table\_36\_tripod\_ai.csv}: each
item is marked satisfied, partly, adapted or not applicable, naming the
section or artefact that addresses it. Items marked \emph{partly} ---
eligibility, outcome timing, fairness, preregistration --- are those the
dataset cannot support.

\end{document}
